\begin{answers}{ch:infra}{計算量の見積もりとインフラ選定}
\item 約 4.8 日で、費用は約 3.1 万ドルです。
$N=10^{10}$、$D=20N=2\times10^{11}$ なので、計算量は $C\approx6ND=1.2\times10^{22}$ FLOPs です。
B200 1 枚の実効の速さは $2{,}250\times10^{12}\times0.4=9\times10^{14}$ FLOP/s です。
式\eqref{eq:infra-gpuhours}より、GPU 時間は $1.2\times10^{22}\div(9\times10^{14}\times3600)\approx3{,}704$ です。
32 枚で割ると約 116 時間、つまり 4.8 日です。
費用は $3{,}704\times8.31\approx30{,}800$ ドルになります。
\item 1 枚には載りません。2 枚なら 52 本です。
BF16 の重みは $2N=1.411\times10^{11}$ バイト、つまり 141.1 GB で、これだけで H200 の 141 GB を超えます。
2 枚（282 GB）なら、重みを引いた残りは $282-141.1=140.9$ GB です。
8,192 トークンの系列 1 本の KV キャッシュは $320\times1024\times8192\approx2.68\times10^{9}$ バイトで、2.68 GB です。
式\eqref{eq:infra-infer-mem}より、持てる本数は $140.9\div2.68\approx52.5$ を切り捨てた 52 本です。
実際には計算の作業用のメモリも要るので、これより少なめに見積もります。
\item 1 GPU 時間が約 14.77 ドルのときです。
\code{GPUS} に \code{"NEXT": (4000, 288, 14.77)} のような 1 行を足します。
1 EFLOP あたりの費用は「価格 $\div$（公称値 $\times$ MFU $\times$ 3600）」です。
MFU が同じなら、価格と公称値の比が等しいときに、費用も等しくなります。
B200 は $8.31\div2{,}250$ なので、求める価格は $8.31\times4{,}000\div2{,}250\approx14.77$ ドルです。
足した GPU で \code{gpu_hours(1e18, "NEXT")} に価格を掛けると、B200 と同じ 2.56 ドルになりました。
この指標にメモリは入らないので、288 GB で分割や通信を減らせるなら、これより高くても選ぶ理由になります。
\item 塊が小さいほど比は大きく、大きいほど小さくなりました（ローカルディスクでの実測）。
小さなファイルを読む時間は、\code{memmap} に対して、1 KiB で約 17 倍、8 KiB で約 3.6 倍、64 KiB で約 2.6 倍でした。
小さなファイル 2,000 個を読む時間は、98 ms、110 ms、167 ms と、あまり変わりません。
\code{memmap} は 5.8 ms、30 ms、66 ms と、大きさにつれて増えました。
\figref{fig:infra-iobench}の右の図で、ファイル 1 個ごとの「開く・読む・閉じる」の往復は、中身の大きさによらず同じ回数だけ起きます。
中身が小さいほど、時間のほとんどがこの往復になります。
中身が大きくなると、どちらの方法でも中身を運ぶ時間が中心になり、比は 1 に近づきます。
ここでは NFS では測っていませんが、往復の待ち時間が長い NFS では、同じ傾向がもっと強く出るはずです。
\item 間隔は約 39 分、失う時間の割合は約 21\% です。
完全なチェックポイントは $12N=12\times4.05\times10^{11}$ バイトで、4,860 GB です。
毎秒 20 GB で書くと、$c=243$ 秒（約 4 分）かかります。
$T_{\mathrm{fail}}=3.1$ 時間 $=11{,}160$ 秒なので、$\tau=\sqrt{2\times243\times11{,}160}\approx2{,}329$ 秒、約 39 分です。
式\eqref{eq:infra-waste}に入れると、$243/2{,}329+2{,}329/(2\times11{,}160)\approx0.104+0.104=0.209$ です。
最もよい間隔では、2 つの項が等しくなります。
本文の 70B の例（同じ速さで 9\%）と比べると、モデルが大きくなった分だけ損失が増えています。
\item 電気代は約 36 万ドルで、借りる費用の約 4\% です。
1 ノードを DGX H100 の 10.2 kW とすると、IT 機器の電力は $31\times10.2=316.2$ kW です。
1 年は 8,760 時間なので、式\eqref{eq:infra-power}より $316.2\times1.3\times0.10\times8{,}760\approx360{,}000$ ドルです。
借りる費用は $248\times3.99\times8{,}760\approx867$ 万ドルです。
割ると約 4.2\% で、本文の例（1 GPU 時間あたり 0.17 ドルと 3.99 ドル）と同じ比になります。
ネットワーク機器やストレージの電力は、ここには入れていません。
\item 例として、3 億パラメータのモデルを $D=20N=60$ 億トークン、H100 1 枚で学習させる計画を考えます。
\code{sheet} の結果は、計算量 $1.08\times10^{19}$ FLOPs、約 8 GPU 時間、約 30 ドル、状態のメモリ 4.8 GB、トークン列 12 GB でした。
最も自信が持てない入力は、多くの場合 MFU です。
本文の検算でも、逆算した MFU は 26〜55\% に散らばりました。
MFU が半分の 20\% なら、時間も費用も 2 倍（約 15 時間、約 61 ドル）になります。
対策は、予算を見積もりの 2〜3 倍にしておく、GPU を増やす、$N$ か $D$ を減らす、のどれかです。
最初の数百ステップで実際の速さを測り、見積もりを直してから本番に進むのが確実です。
\end{answers}
